\section{Absorption models and the method of characteristics}
\label{sec:moc}

When one compartment ruptures, its contents enter a shared medium before they are
taken up by the compartments that remain.  Treating this transfer as instantaneous
suppresses both the absorption time and any loss that occurs after uptake.  The
models below retain an explicit absorption period.  They follow \(X_t\), the number
of particles inside one receiving compartment, and \(Y_t\), the number still
outside, from the initial state
\begin{equation}
  X_0=0,
  \qquad Y_0=N.
  \label{eq:abs-initial-state}
\end{equation}
The three models are deliberately laddered.  Pure absorption establishes the
bivariate PGF.  Interior death adds a third fate but leaves a finite state space.
Interior birth then makes the state space infinite and turns the PGF equation into
the non-trivial characteristic problem required by later compartment models.

For all three models write
\begin{equation}
  p_{ij}(t)=\Prb(X_t=i,Y_t=j),
  \qquad
  G(x,y,t)=\sum_{i,j\geq0}p_{ij}(t)x^iy^j
  =\E[x^{X_t}y^{Y_t}].
  \label{pgfDef}
\end{equation}
Probabilities with a negative index are understood to be zero.

\subsection{Absorption only}
\label{sec:absonly}

Suppose each exterior particle is absorbed independently at rate \(\alpha>0\).
From state \((i,j)\),
\begin{align}
  \Prb(\Delta X=1,\Delta Y=-1\mid X_t=i,Y_t=j)
  &=\alpha j\,\Delta t+o(\Delta t),
  \label{eq:abs-short-time}\\
  \Prb(\Delta X=0,\Delta Y=0\mid X_t=i,Y_t=j)
  &=1-\alpha j\,\Delta t+o(\Delta t).
  \nonumber
\end{align}
The forward equations are therefore
\begin{equation}
  \frac{\dd p_{ij}}{\dd t}
  =\alpha(j+1)p_{i-1,j+1}-\alpha j p_{ij}.
  \label{eq:abs-master}
\end{equation}
Multiplying by \(x^iy^j\), summing and shifting the incoming indices gives
\begin{equation}
  \frac{\partial G}{\partial t}
  =\alpha(x-y)\frac{\partial G}{\partial y},
  \qquad G(x,y,0)=y^N.
  \label{PDE1}
\end{equation}

The probabilistic solution is immediate.  One initial particle is still outside
at time \(t\) with probability \(\e^{-\alpha t}\) and has been absorbed with the
complementary probability.  Its PGF is therefore
\[
  g_1(x,y,t)=\bigl(1-\e^{-\alpha t}\bigr)x+\e^{-\alpha t}y.
\]
The \(N\) initial particles act independently, so
\begin{equation}
  G(x,y,t)=
  \left[
    \bigl(1-\e^{-\alpha t}\bigr)x+\e^{-\alpha t}y
  \right]^N.
  \label{eq:Gabsonly}
\end{equation}
It follows that
\(X_t\sim\operatorname{Binomial}(N,1-\e^{-\alpha t})\) and
\(Y_t=N-X_t\).  Substitution verifies that \eqref{eq:Gabsonly} also solves the
PDE \eqref{PDE1}.

\subsection{Absorption with interior death}
\label{sec:absdeath}

Now let each absorbed particle die at rate \(\mu\geq0\).  The additional
transition \((i,j)\to(i-1,j)\) has rate \(\mu i\), so
\begin{equation}
  \frac{\dd p_{ij}}{\dd t}
  =\alpha(j+1)p_{i-1,j+1}
  +\mu(i+1)p_{i+1,j}
  -(\alpha j+\mu i)p_{ij}.
  \label{eq:absdeath-master}
\end{equation}
The bivariate PGF satisfies
\begin{equation}
  \frac{\partial G}{\partial t}
  =\mu(1-x)\frac{\partial G}{\partial x}
  +\alpha(x-y)\frac{\partial G}{\partial y},
  \qquad G(x,y,0)=y^N.
  \label{PDE2}
\end{equation}

For a single initial particle there are three possible states at time \(t\): still
outside, alive inside, or dead.  Their probabilities are
\begin{align}
  q_{\mathrm o}(t)&=\e^{-\alpha t},
  \label{eq:qoutside}\\
  q_{\mathrm i}(t)&=
  \begin{cases}
    \displaystyle
    \frac{\alpha}{\mu-\alpha}
    \bigl(\e^{-\alpha t}-\e^{-\mu t}\bigr),
    &\mu\neq\alpha,\\[9pt]
    \alpha t\e^{-\alpha t},&\mu=\alpha,
  \end{cases}
  \label{eq:qinside}\\
  q_{\mathrm d}(t)&=1-q_{\mathrm o}(t)-q_{\mathrm i}(t).
  \label{eq:qdead}
\end{align}
The second line is the convolution of absorption at time \(s\) with survival
inside for the remaining time \(t-s\); its equal-rate expression is the regular
limit omitted by a quotient-only formula.  Independence now gives the trinomial
PGF
\begin{equation}
  G(x,y,t)
  =\bigl[q_{\mathrm d}(t)+q_{\mathrm i}(t)x+q_{\mathrm o}(t)y\bigr]^N.
  \label{eq:Gabsdeath}
\end{equation}

\begin{figure}[htbp]
  \centering
  \begin{subfigure}[t]{0.47\textwidth}
    \centering
    \includegraphics[width=\textwidth]{abs1.pdf}
    \caption{Absorption only.}
  \end{subfigure}
  \hfill
  \begin{subfigure}[t]{0.47\textwidth}
    \centering
    \includegraphics[width=\textwidth]{abs2.pdf}
    \caption{Absorption with interior death.}
  \end{subfigure}
  \caption{Single-particle state probabilities.  Pure absorption transfers mass
  from the exterior to the interior.  Adding death makes the interior probability
  rise and then fall as mass is transferred to the absorbing dead state.}
  \label{fig:absorption-probabilities}
\end{figure}

\subsection{Characteristics on the finite-state model}
\label{sec:mocabsdeath}

The direct solution \eqref{eq:Gabsdeath} provides a controlled test of the method
of characteristics.  Write \eqref{PDE2} as
\[
  -G_t+\mu(1-x)G_x+\alpha(x-y)G_y=0
\]
and introduce characteristic coordinates \((\xi,\eta,\tau)\) with initial surface
\begin{equation}
  x(0)=\xi,
  \qquad y(0)=\eta,
  \qquad t(0)=0.
  \label{eq:characteristic-initial-surface}
\end{equation}
Along a characteristic,
\begin{equation}
  \frac{\dd G}{\dd\tau}=0,
  \qquad
  \frac{\dd t}{\dd\tau}=-1,
  \qquad
  \frac{\dd x}{\dd\tau}=\mu(1-x),
  \qquad
  \frac{\dd y}{\dd\tau}=\alpha(x-y).
  \label{eq:absdeath-characteristics}
\end{equation}
The first two equations give \(G=\eta^N\) and \(t=-\tau\).  The third gives
\begin{equation}
  1-x=(1-\xi)\e^{-\mu\tau}.
  \label{eq:absdeath-x-characteristic}
\end{equation}
Substitution into the \(y\)-equation, followed by one integrating factor, yields
for \(\alpha\neq\mu\)
\begin{equation}
  y=1-\frac{\alpha}{\alpha-\mu}(1-\xi)\e^{-\mu\tau}
  +\left[
      \eta+\frac{\alpha}{\alpha-\mu}(1-\xi)-1
    \right]\e^{-\alpha\tau}.
  \label{eq:absdeath-y-characteristic}
\end{equation}
Eliminating \(\xi\) and \(\tau\) and solving for the initial coordinate
\(\eta\) gives
\[
  \eta=q_{\mathrm d}(t)+q_{\mathrm i}(t)x+q_{\mathrm o}(t)y.
\]
Since \(G=\eta^N\), the characteristic calculation reproduces
\eqref{eq:Gabsdeath}.  The equal-rate case follows by taking the limit, giving
\eqref{eq:qinside}; the machinery is therefore fixed before the infinite-state
model is introduced.

\subsection{Absorption with interior birth and death}
\label{sec:absbirthdeath}

Finally let each interior particle give birth at rate \(\beta\geq0\) and die at
rate \(\mu\geq0\).  The three transitions are
\begin{align*}
  (i,j)&\longrightarrow(i+1,j-1) &&\text{at rate }\alpha j,\\
  (i,j)&\longrightarrow(i+1,j)   &&\text{at rate }\beta i,\\
  (i,j)&\longrightarrow(i-1,j)   &&\text{at rate }\mu i.
\end{align*}
Their forward equation is
\begin{multline}
  \frac{\dd p_{ij}}{\dd t}
  =\alpha(j+1)p_{i-1,j+1}
  +\mu(i+1)p_{i+1,j}
  +\beta(i-1)p_{i-1,j}\\
  -(\alpha j+\mu i+\beta i)p_{ij},
  \label{eq:abd-master}
\end{multline}
and the PGF satisfies
\begin{equation}
  G_t=
  \bigl[\mu-(\mu+\beta)x+\beta x^2\bigr]G_x
  +\alpha(x-y)G_y,
  \qquad G(x,y,0)=y^N.
  \label{PDE3}
\end{equation}

Birth makes infinitely many interior states reachable, but independence has not
disappeared.  Descendants of distinct exterior particles remain independent, so
\(G=h^N\), where \(h\) is the PGF generated by one exterior particle.  What has
changed is that \(h\) is no longer a finite polynomial.

There is a branch-safe representation of \(h\) that is valid for every
non-negative choice of the rates.  Let \(g(x,u)\) be the PGF of a linear
birth--death population of age \(u\), started from one interior particle.  From
\cref{sec:linear-bd},
\begin{equation}
  g(x,u)=
  \begin{cases}
  \displaystyle
  \frac{\mu(1-x)-(\mu-\beta x)\e^{-(\beta-\mu)u}}
       {\beta(1-x)-(\mu-\beta x)\e^{-(\beta-\mu)u}},
  &\beta\neq\mu,\\[14pt]
  \displaystyle
  1-\frac{1-x}{1+\beta u(1-x)},
  &\beta=\mu.
  \end{cases}
  \label{eq:interior-bd-pgf}
\end{equation}
If the exterior particle has not been absorbed by time \(t\), it contributes
\(y\); if absorption occurs at time \(s\leq t\), it initiates an interior process
of age \(t-s\).  Conditioning on the absorption time therefore gives
\begin{equation}
  h(x,y,t)
  =\e^{-\alpha t}y
  +\int_0^t\alpha\e^{-\alpha s}g(x,t-s)\,\dd s,
  \qquad
  G(x,y,t)=h(x,y,t)^N.
  \label{eq:abd-integral-solution}
\end{equation}
This finite integral is the preferred global formula.  It remains regular when
\(\beta=\mu\), when \(\beta=0\), and at the resonant rate combinations that make
individual terms in a special-function expansion singular.

\subsection{Characteristic and hypergeometric form}

The characteristic system for \eqref{PDE3} has the same \(G\), \(t\) and \(y\)
equations as \eqref{eq:absdeath-characteristics}, while
\begin{equation}
  \frac{\dd x}{\dd\tau}
  =\mu-(\mu+\beta)x+\beta x^2.
  \label{eq:abd-x-characteristic}
\end{equation}
Assume for the moment that \(\beta>0\) and
\(b\defeq\beta-\mu\neq0\).  The right-hand side factorises as
\begin{equation}
  \beta(x-x_+)(x-x_-),
  \qquad x_+=1,
  \qquad x_-=\frac{\mu}{\beta},
  \label{eq:abd-fixed-points}
\end{equation}
and separation gives
\begin{equation}
  \frac{x-x_+}{x-x_-}
  =\frac{\xi-x_+}{\xi-x_-}\e^{b\tau}.
  \label{eq:xratio}
\end{equation}
Write
\[
  Z=\frac{\xi-x_+}{\xi-x_-},
  \qquad
  \zeta(\tau)=Z\e^{b\tau}.
\]
Solving \eqref{eq:xratio} for \(x\) gives
\begin{equation}
  x(\tau)=x_++(x_+-x_-)
  \frac{\zeta(\tau)}{1-\zeta(\tau)}.
  \label{eq:xgeom}
\end{equation}
The remaining characteristic is
\begin{equation}
  \frac{\dd}{\dd\tau}\bigl(\e^{\alpha\tau}y\bigr)
  =\alpha\e^{\alpha\tau}x(\tau).
  \label{eq:yint}
\end{equation}
Where \(\lvert\zeta\rvert<1\), expand the quotient geometrically and integrate
term by term.  With
\begin{equation}
  \delta=\frac{\alpha}{b},
  \qquad
  \Psi(z)=1+\frac{\alpha}{\beta}
  \sum_{n=1}^{\infty}\frac{z^n}{n+\delta},
  \label{eq:Psidef}
\end{equation}
one obtains
\(
  \alpha\int\e^{\alpha\tau}x(\tau)\dd\tau
  =\e^{\alpha\tau}\Psi(\zeta(\tau))
\)
up to a constant.  At the endpoint \((x,y,t)\), set
\begin{equation}
  W=\frac{x-1}{x-\mu/\beta}.
  \label{eq:Wdef}
\end{equation}
Because \(t=-\tau\), the initial characteristic ratio is
\(Z=W\e^{bt}\).  Solving \eqref{eq:yint} for the initial coordinate \(\eta\)
and using \(G=\eta^N\) gives the closed expression
\begin{equation}
  G(x,y,t)=
  \left\{
    \e^{-\alpha t}\bigl[y-\Psi(W)\bigr]
    +\Psi\bigl(W\e^{bt}\bigr)
  \right\}^{N}.
  \label{eq:Gabd}
\end{equation}

Equation \eqref{eq:Gabd} is first obtained in the joint convergence region
\begin{equation}
  |W|<1,
  \qquad
  |W|\e^{bt}<1,
  \label{eq:Psi-series-domain}
\end{equation}
with \(n+\delta\neq0\) for every \(n\geq1\).  Within that region it agrees with
the integral solution \eqref{eq:abd-integral-solution}, and the identity supplies
its analytic continuation wherever a consistent branch is chosen.  For numerical
work outside \eqref{eq:Psi-series-domain}, especially in the supercritical regime
at \(x\leq\mu/\beta\), the finite integral is safer: \(W\) may cross the standard
hypergeometric branch cut even though the physical PGF remains real and analytic.

The series is a Gauss hypergeometric function.  Since
\begin{equation}
  {}_2F_1(1,\delta;\delta+1;z)
  =\sum_{n=0}^{\infty}\frac{\delta}{n+\delta}z^n,
  \label{eq:2F1simple}
\end{equation}
we may write
\begin{equation}
  \Psi(z)
  =1+\frac{\beta-\mu}{\beta}
  \left[{}_2F_1(1,\delta;\delta+1;z)-1\right].
  \label{eq:Psihyp}
\end{equation}
This form is compact, but it does not remove the need to specify analytic
continuation.

Three checks are independent of that representation.  At \(t=0\), the two
\(\Psi\) terms in \eqref{eq:Gabd} cancel and \(G=y^N\).  At \(x=y=1\),
\(W=0\) and \(\Psi(0)=1\), giving \(G=1\).  Finally,
\begin{equation}
  G(1,y,t)=
  \bigl[1-\e^{-\alpha t}+\e^{-\alpha t}y\bigr]^N,
  \label{eq:abd-exterior-marginal}
\end{equation}
so \(Y_t\sim\operatorname{Binomial}(N,\e^{-\alpha t})\), as it must: interior
birth and death cannot alter whether an exterior founder has been absorbed.

Two degenerate cases now have a precise status.  If \(\beta=\mu\), the roots in
\eqref{eq:abd-fixed-points} coalesce and \(b\), \(\delta\) and \(W\) cannot be used;
the critical PGF in \eqref{eq:interior-bd-pgf} inserted into
\eqref{eq:abd-integral-solution} is the exact replacement.  If
\begin{equation}
  \alpha=n(\mu-\beta)
  \qquad\text{for some }n\geq1,
  \label{eq:abd-resonance}
\end{equation}
then a denominator \(n+\delta\) in \eqref{eq:Psidef} vanishes.  The singularity
is removable in the difference of \(\Psi\) terms in \eqref{eq:Gabd}; its limit
contains a factor of \(t\) multiplying the resonant exponential.  Again,
\eqref{eq:abd-integral-solution} evaluates the solution directly without a
parameter limit.

The progression from \eqref{PDE1} to \eqref{PDE3} is the part that later models
reuse.  State transitions determine the master equation; the PGF compresses it;
characteristics expose the closed form where one exists; and the probabilistic
integral keeps the assumptions and exceptional parameter values explicit.

\FloatBarrier
